Tato sekce shrnuje postupy pro použití interaktivních AI nástrojů ke podpoře studentů s odlišným jazykovým zázemím a seznámení s neurodiverzitou. Cílem jsou rychlé, opakovatelné šablony a kontrolní body pro výuku i online aktivity.

Praktické principy
\begin{itemize}
  \item Lokalizace a kulturní adaptace: neprovádějte mechanický překlad—ověřte příklady vůči cílové skupině.
  \item Úprava tempa a složitosti: připravte minimálně dvě verze (zjednodušená / standardní) a možnost krokového rozkladu.
  \item Multimodalita: kombinujte text, audio a vizuály; jednoduchý diagram nebo piktogram často zefektivní pochopení. Pro rychlou tvorbu vizuálů použijte Microsoft Designer \cite{kb_ai_101_tipu_pdf_04e4a115_page_51_1}.
  \item Důstojnost a autonomie: AI navrhuje podporu, student má vždy možnost ji přijmout nebo odmítnout.
\end{itemize}

Konkrétní nástroje (ilustrativně)
\begin{description}
  \item[Minecraft Education:] vhodný pro vizuální, interaktivní lekce a připravené světy (např. Hour of Code) pro studenty s omezenou znalostí jazyka \cite{kb_ai_101_tipu_pdf_04e4a115_page_15_2}.
  \item[Microsoft Designer:] rychlé generování piktogramů, diagramů a krokových návodů pro doplnění textu \cite{kb_ai_101_tipu_pdf_04e4a115_page_51_1}.
  \item[OneNote / math.microsoft.com:] převod rukopisu na digitální rovnice a analýza postupu řešení (užitečné v matematice) \cite{kb_ai_101_tipu_pdf_04e4a115_page_8_1}.
\end{description}

Pracovní workflow (pro jedno cvičení)
\begin{enumerate}
  \item Definujte cíl podpory: jazyková pomoc, kognitivní rozklad nebo multimodální materiál.
  \item Připravte základní zadání a dvě upravené verze: zjednodušenou a rozšířenou.
  \item Vytvořte stručný vizuál (diagram/piktogram) a sadu konzistentních ikon (Designer).
  \item Volitelně vložte aktivitu do interaktivního prostředí (Minecraft Education).
  \item Požádejte o nahrání postupu (foto rukopisu, audio/video) a využijte převod rukopisu k analýze.
  \item Poskytněte krátkou personalizovanou zpětnou vazbu a alternativní hint (obrázek nebo 30–60 s audio).
\end{enumerate}

Šablony promptů (kopírovatelně)
\begin{itemize}
  \item Lokalizace / jednoduchý překlad:
    \begin{quote}
      
\begin{PromptBlock}

Přelož následující text do [JAZYK] a vytvoř zjednodušenou verzi pro studenty s nižší jazykovou úrovní (max. 5 vět). Poté navrhni krátký obrázkový návod na 3 krocích.
\end{PromptBlock}

    \end{quote}
  \item Víceúrovňové vysvětlení:
    \begin{quote}
      
\begin{PromptBlock}

Vysvětli [POJEM] ve třech úrovních: 1) velmi jednoduché (1–2 věty), 2) střední s příkladem (3–4 věty), 3) detailní s kroky.
\end{PromptBlock}

    \end{quote}
  \item Generování multimodálního hintu:
    \begin{quote}
      
\begin{PromptBlock}

Pro úlohu [stručný popis] vygeneruj: a) krátké textové vodítko (max. 30 slov), b) návrh diagramu s 4 popisky, c) audio skript pro ~45s.
\end{PromptBlock}

    \end{quote}
\end{itemize}

Checklist pro nasazení
\begin{enumerate}
  \item Zkontrolujte jazyková nastavení a dostupnost překladu.
  \item Připravte minimálně dvě úrovně obtížnosti a multimodální verze materiálu.
  \item Zajistěte možnost nahrávání postupu a workflow pro jeho zpracování (OneNote apod.).
  \item Použijte vizuální pomůcky tam, kde text nestačí.
  \item Informujte studenty o roli AI a podmínkách uchovávání materiálů (souhlas/retence).
\end{enumerate}

Evaluace účinnosti
\begin{itemize}
  \item Po pilotu sbírejte krátké ankety o srozumitelnosti, přístupnosti a vlivu na učení.
  \item Měřte využití modalit (audio vs. vizuál) a upravujte materiály podle výsledků.
  \item Pravidelně aktualizujte šablony a vizuální styly dle zpětné vazby.
\end{itemize}

Poznámka o zdrojích
\begin{itemize}
  \item Minecraft Education nabízí připravené světy pro demonstrace a Hour of Code \cite{kb_ai_101_tipu_pdf_04e4a115_page_15_2}.
  \item Microsoft Designer usnadňuje rychlou tvorbu vizuálů v různých stylech \cite{kb_ai_101_tipu_pdf_04e4a115_page_51_1}.
  \item Nástroje Microsoft pro převod rukopisu podporují práci s postupy v matematice \cite{kb_ai_101_tipu_pdf_04e4a115_page_8_1}.
\end{itemize}